\documentclass{scrreprt}

\usepackage[croatian]{babel}
\usepackage[utf8]{inputenc}
\usepackage[T1]{fontenc}
\usepackage{hyperref}
\usepackage{graphicx}
\usepackage{tikz}
\usepackage{pgfplots}
\usepackage{float}
\usepackage[nottoc,numbib]{tocbibind}
\usepackage{color}
\usepackage{multicol}
\usepackage[export]{adjustbox}

\setlength{\parskip}{\bigskipamount}
\setlength{\parindent}{0pt}

\graphicspath{ {./images/} }

\begin{document}

\titlehead{Sveučilište u Zagrebu\\Filozofski fakultet\\Odsijek za informacijske
i komunikacijske znanosti\\Akademska godina 2013/14.}
\title{Upotreba web aplikacije za kvizove u nastavi osnovnih i srednjih škola}
\author{Studenti: Janko i Matija Marohnić\\Mentor: doc.dr.sc. Kristina Kocijan}
\date{Zagreb, 2014.}

\maketitle

\pagebreak

Ovaj rad izrađen je na Odsijeku za informacijske i komunikacijske znanosti na
Filozofskom fakultetu u Zagrebu pod vodstvom prof. dr. sc. Kristine Kocijan i
predan je na natječaj za dodjelu Rektorove nagrade u akademskoj godini
2013/2014.

\pagebreak

\tableofcontents

\chapter{Uvod}

\emph{Kvizovi} su aplikacija za riješavanje kvizova koja je prilagođena osnovnim
i srednjim školama tako što se u korijenu dijeli na dvije uloge -- učenike i
profesore.

Postoje razni kreativni načini da se upotrijebi informacijska tehnologija u
nastavi, ali to je obično prilično jednostrano, odnosno, samo profesori koriste
tehnologiju dok ih učenici gledaju. Smatramo da se trebaju razvijati puno
interaktivnije i zabavnije metode podučavanja.

U ovom smo istraživanju testirali našu aplikaciju za kvizove na nekoliko
osnovnih i srednjih škola. Cilj ovog istraživanja bio je saznati:

\begin{itemize}
  \item pomaže li takav način učenja učenicima da bolje savladaju gradivo
  \item čini li aplikacija profesorima podučavanje zanimljivijim
\end{itemize}

\chapter{Hipoteza}

Naša je pretpostavka da učenici koji su učili riješavajući kvizove pomoću
aplikacije bolje usvajaju gradivo od onih koji nisu i da je profesorima na taj
način zanimljivije provoditi ispitivanje gradiva.

\chapter{Sredstva i metode}

Istraživanje smo proveli tako što smo izradili aplikaciju za kvizove i dali
određenom uzorku osnovnih i srednjih škola na testiranje. Podatke o kvizovima
skupljali smo 2 godine. U svakom razredu su učenici bili podijeljeni na skupinu
koja ne riješava kvizove i na onu koja riješava, na taj način smo mogli
provjeravati učinkovitost aplikacije kada bi učenici imali test znanja. Imamo
puno podataka na raspolaganju, ali smo za svaki slučaj proveli još i ankete kako
bismo dobili podatke koje pasivno ne bismo mogli dobiti.

Imamo blog na kojemu obaviještavamo korisnike o ažuriranju aplikacije. S obzirom
da je dio ciljane publike starija populacija koja nije nužno iskusna s
aplikacijama, dodali smo i vodič kroz aplikaciju koji je dostupan u bilo kojem
trenutku, kako bi se korisnici lakše snašli. Također, u slučaju da ima nekih
problema, korisnici nas uvijek mogu kontaktirati putem e-mail adresa koje se
nalaze u kontakt sekciji koju smo stavili u glavnu navigaciju. Korisnici također
mogu naknadno izmjenjivati svoje podatke u korisničkom računu.

Korisnik započinje korištenje aplikacije tako da odabere svoju ulogu u školi,
stoga ovisno o ulozi aplikacija pruža drugačiju funkcionalnost.

\begin{figure}[H]
  \includegraphics[width=\textwidth, clip=true, trim=0 7cm 0 0, fbox]{home}
  \caption{Početna stranica}
\end{figure}

\section{Škola}

\emph{Škola} je uloga koja predstavlja profesora i ona je administrator kvizova.

Nakon odabira te uloge korisnik se može prijaviti ili registrirati ako još nema
korisnički račun. Registracija se sastoji od ispunjavanja jednostavnog formulara
pomoću kojega skupljamo informacije o korisnicima koje možemo iskoristiti kako
bismo poboljšali njihovo iskustvo i kako bismo mogli raditi istraživanja.
Zanimljivo polje u formularu za registraciju je \emph{Tajni ključ}, koji je
potreban za registraciju učenicima te škole.

Nakon prijave profesore dočeka lista kvizova koje su napravili,

\begin{figure}[H]
  \includegraphics[width=\textwidth, clip=true, trim=0 10cm 0 0, fbox]{school/quizzes}
  \caption{Škole -- popis kvizova}
\end{figure}

gdje mogu izmjenjivati postojeće kvizove i sastavljati nove. Nakon što je
profesor zadovoljan s kvizom, može ga učiniti aktivnim, odnosno vidljivim
učenicima. Izmjenjivanje kvizova podijeljeno je na izmjenu metapodataka kviza i
na izmjenu pitanja kviza.

\begin{figure}[H]
  \includegraphics[width=\textwidth, clip=true, trim=0 53cm 0 0, fbox]{school/quiz}
  \caption{Škole -- popis pitanja u kvizu}
\end{figure}

Postoji 4 vrsta pitanja:

\begin{itemize}
  \item Točno/netočno
  \item Ponuđeni odgovori
  \item Asocijacija
  \item Upiši točan odgovor
\end{itemize}

Uz svako se pitanje može pridružiti pomoć, koja će se prikazati učenicima dok
riješavaju pitanje.

Škole mogu pregledavati informacije o odigrane kvizove, kada su igrani, tko ih
je igrao, koji su bili rezultati itd.

\begin{figure}[H]
  \includegraphics[width=\textwidth, clip=true, trim=0 10cm 0 0, fbox]{school/played_quizzes}
  \caption{Škole -- popis odigranih kvizova}
\end{figure}

\begin{figure}[H]
  \includegraphics[width=\textwidth, clip=true, trim=0 80cm 0 0, fbox]{school/played_quiz}
  \caption{Odigrani kviz}
\end{figure}

Mogu se pregledavati i učenici, koliko su kvizova odigrali, koji su to kvizovi i
sl.

\begin{figure}[H]
  \includegraphics[width=\textwidth, clip=true, trim=0 15cm 0 0, fbox]{school/students}
  \caption{Škole -- popis učenika}
\end{figure}

\section{Učenik}

\emph{Učenik} je uloga koja riješava kvizove koje je napravila njihova škola.
Kao i kod škole, učenik se može prijaviti ili registrirati, ako već nema
korisnički račun. Pri registraciji učenik treba napisati tajni ključ koji mu je
njegova škola dala, u protivnom se ne može registrirati. Na taj način
spriječavamo da se bilo tko registrira kao učenik.

Nakon prijave ili registracije, korisnika dočeka lista kvizova koji su dostupni
za riješavanje. Kviz je moguće igrati sam ili u paru, u drugom slučaju se drugi
igrač također treba prijaviti.

\begin{figure}[H]
  \includegraphics[width=\textwidth, clip=true, trim=0 8.5cm 0 0, fbox]{student/quizzes}
  \caption{Učenici -- popis kvizova}
\end{figure}

Nakon što učenik započne kviz, prikazuje mu se jedno po jedno pitanje na koja
treba odgovoriti. Kako bi igra bila što poštenija, pitanja su obično vremenski
ograničena, tako da se učenik ne stigne previše konzultirati s vanjskim
izvorima.

\begin{figure}[H]
  \includegraphics[width=\textwidth, clip=true, trim=0 7cm 0 0, fbox]{student/boolean_question}
  \caption{Učenici -- točno/netočno}
\end{figure}

\begin{figure}[H]
  \includegraphics[width=\textwidth, clip=true, trim=0 7cm 0 0, fbox]{student/boolean_question_correct}
  \caption{Učenici -- točno/netočno s povratnom informacijom}
\end{figure}

Nakon što učenik odgovori na sva pitanja, ispisuju se rezultati i učenik dobiva
određenu ``titulu'' s obzirom na njegov rezultat. Titule su osmišljene tako da
se učenika uvijek pohvaljuje, čak i ako je imao loš rezultat, tako da se učenik
dobro osjeća i da ga se potiče da igra i dalje.

% \chapter{Tablice i grafovi}

% \begin{table}
%   \caption{Anketa za škole}

%   \begin{multicols}{2}
%     \small{
%       \begin{enumerate}
%         \item Spol, vrsta škole, naziv škole i zanimanje
%         \item Koliko dugo radite u školi?
%         \item Koje predmete predajete?
%         \item Koristite li računalo u nastavi?
%         \item Ako da, koliko često?
%           \begin{itemize}
%             \item Svaki sat
%             \item Jednom tjedno
%             \item Jednom mjesečno
%             \item Jednom u polugodištu
%             \item Jednom u cijeloj godini
%           \end{itemize}
%         \item Koristite li računalo za pripremu nastave?
%         \item Ako da, koliko često?
%           \begin{itemize}
%             \item Sve pripreme radim na računalu
%             \item Tjedne pripreme
%             \item Mjesečne pripreme
%             \item Plugodišnje pripreme
%             \item Godišnje pripreme
%             \item Samo za unos ocjena na kraju godine
%           \end{itemize}
%         \item Koliko ste dugo koristili aplikaciju Kvizovi?
%         \item Jeste li imali poteškoća prilikom korištenja aplikacije?
%         \item Ako da, koje ste poteškoće imali?
%         \item Biste li aplikaciju koristili i dalje u svom radu?
%         \item Biste li aplikaciju preporučili svojim kolegama iz istog nastavnog
%           predmeta?
%         \item Biste li aplikaciju preporučili kolegama iz drugih nastavnih
%           predmeta
%           \begin{itemize}
%             \item Ne, ne bih je nikome preporučio/la
%             \item Samo bih je preporučio/la kolegama iz hrvatskog jezika
%             \item Preporučio/la bih je svim jezičnim predmetnim nastavnicima
%             \item Preporučio/la bih ga svim nastavnicima prirodnih znanosti
%             \item Preporučio/la bih ga svim nastavnicima društvenih znanosti
%             \item Preporučio/la bih ga svim nastavnicima tehničkih znanosti
%           \end{itemize}
%         \item Kako biste ocijenili znanje iz lektire one djece koja su koristila
%           aplikaciju (nakon korištenja)?
%           \begin{itemize}
%             \item Nisam primjetio/la da su napredovali zbog korištenja
%               aplikacije
%             \item Za nijansu su bili bolji nego dok nisu koristili aplikaciju
%             \item Značajno im se poboljšalo poznavanje gradiva iz lektire
%             \item Uz bolje poznavanje gradiva, primjetio/la sam i veću
%               zainteresiranost djece za lektire i njihovo bolje sudjelovanje u
%               nastavi lektire
%           \end{itemize}
%         \item Koliko mislite da je aplikacija:
%           \begin{itemize}
%             \item Više odmogla nego pomogla učenicima
%             \item Nije doprinijela boljem poznavanju gradiva
%             \item Djelomično bila od pomoći
%             \item Bila od iznimne pomoći kod stvaranja većeg interesa za lektire
%               i poznavanje gradiva
%           \end{itemize}
%         \item U usporedbi s djecom koja nisu koristila aplikaciju Lektire:
%           \begin{itemize}
%             \item djeca koja nisu koristila aplikaciju bolje su znala gradivo od
%               djece koja \textbf{jesu} koristila aplikaciju.
%             \item svi su pokazali podjednako znanje.
%             \item djeca koja su koristila aplikaciju bolje su znala od djece
%               koja \textbf{nisu} koristila aplikaciju.
%           \end{itemize}
%         \item Molimo vas da ocijenite sljedeće aspekte aplikacije:
%           \begin{itemize}
%             \item Obaviještenost o promijenama
%             \item Snalaženje u sučelju
%             \item Izgled aplikacije
%           \end{itemize}
%         \item Što vam se osobito svidjelo u aplikaciji?
%         \item Koju mogućnost biste htjeli da aplikacija u budućnosti ima, a
%           trenutno ju nema?
%       \end{enumerate}
%     }
%   \end{multicols}

%   \label{survschools}
% \end{table}

% \begin{table}
%   \caption{Anketa za učenike}

%   \begin{multicols}{2}
%     \small{
%       \begin{enumerate}
%         \item Spol, vrsta škole, naziv škole, razred
%         \item Molimo vas da ocijenite sljedeće aspekte aplikacije:
%           \begin{itemize}
%             \item Obaviještenost o promijenama
%             \item Snalaženje u sučelju
%             \item Izgled aplikacije
%           \end{itemize}
%         \item Smatrate li ovaj način učenja zanimljivim?
%           \begin{itemize}
%             \item Da
%             \item Ne znam
%             \item Ne
%           \end{itemize}
%         \item Smatrate li ovakav način učenja učinkovitim? Odnosno, smatrate li
%           da ste puno naučili u odnosu na uloženi trud?
%           \begin{itemize}
%             \item Da
%             \item Ne znam
%             \item Ne
%           \end{itemize}
%         \item Koliko često koristite aplikaciju?
%           \begin{itemize}
%             \item Samo jednom
%             \item Jednom mjesečno
%             \item Nekoliko puta mjesečno
%             \item Jednom tjedno
%             \item Nekoliko puta tjedno
%             \item Jednom dnevno
%             \item Više puta dnevno
%           \end{itemize}
%         \item Jeste li imali poteškoće prilikom korištenja aplikacije?
%         \item Ako da, koje ste poteškoće imali?
%         \item Je li vam aplikacija pomogla kod:
%           \begin{itemize}
%             \item Razumijevanja lektire
%             \item Zainteresiranosti za lektiru
%             \item Osjećaja samouvjerenosti kod pravog ispitivanja na satu
%           \end{itemize}
%         \item Otkud ste pristupali aplikaciji?
%           \begin{itemize}
%             \item Od kuće s vlastitog računala
%             \item Od kuće s roditeljskog računala
%             \item Iz školske učionice hrvatskog jezika
%             \item Iz školske knjižnice
%             \item Iz gradske knjižnice
%             \item Kod prijatelja
%           \end{itemize}
%         \item Kad ste koristili aplikaciju?
%           \begin{itemize}
%             \item Prije čitanja lektire koja se testira
%             \item Nakon čitanja lektire koja se testira
%             \item Za vrijeme čitanja lektire koja se testira
%           \end{itemize}
%         \item Nakon pristupa aplikaciji, vaš ocjena iz lektire
%           \begin{itemize}
%             \item se poboljšala
%             \item ostala je ista
%             \item se smanjila
%           \end{itemize}
%         \item Na koji vas način profesor/ica ispituje lektiru na satu?
%           \begin{itemize}
%             \item Isto kao i u aplikaciji
%             \item Djelomično kao i u aplikaciji
%             \item Posve drugačije nego u aplikaciji
%           \end{itemize}
%         \item U usporedbi s onim što pita profesor, pitanja u aplikaciji su:
%           \begin{itemize}
%             \item puno lakša
%             \item podjednako laka
%             \item podjednako teška
%             \item puno teža
%           \end{itemize}
%         \item Biste li voljeli da se ovakav način učenja organizira i u drugim
%           predmetima? Ako da, u kojim?
%         \item Što mislite o natjecanju u znanju iz lektire s drugim razredima i
%           školama u Hrvatskoj?
%           \begin{itemize}
%             \item Dodatno bi me poticalo u čitanju lektira
%             \item Ne bi me dodatno poticalo
%           \end{itemize}
%         \item Ako bi vas dodatno poticalo, na koji način biste se najradije
%           htjeli natjecati?
%           \begin{itemize}
%             \item Pojedinačno
%             \item Kao razred
%             \item Kao škola
%           \end{itemize}
%         \item Kako vam se svidjela aplikacija tijekom vremena?
%           \begin{itemize}
%             \item Na početku
%             \item U sredini
%             \item Pri kraju
%           \end{itemize}
%         \item Što ste očekivali od aplikacije?
%           \begin{itemize}
%             \item Ništa
%             \item Ništa, samo zabavu
%             \item Bolje usvajanje znanja iz lektire
%             \item Bolje usvajanje znanja i povećan interes
%           \end{itemize}
%         \item Što ste dobili od aplikacije?
%           \begin{itemize}
%             \item Ništa
%             \item Ništa, samo zabavu
%             \item Bolje usvajanje znanja iz lektire
%             \item Bolje usvajanje znanja i povećan interes
%           \end{itemize}
%         \item Što vam se osobito svidjelo u aplikaciji?
%         \item Koju mogućnost biste htjeli da aplikacija u budućnosti ima, a trenutno
%           ju nema?
%       \end{enumerate}
%     }
%   \end{multicols}

%   \label{survstudents}
% \end{table}

\chapter{Rezultati}

Ovo istraživanje pokazalo je da je većina profesora primjetila napredak u
poznavanju gradiva kod svojih učenika, kao i veću zainteresiranost za gradivo.

\begin{figure}[H]
  \includegraphics[width=\textwidth, clip=true, trim=0 2.5cm 0 0]{advance}
  \caption{Napredak učenika}
\end{figure}

\begin{figure}[H]
  \includegraphics[width=\textwidth, clip=true, trim=0 2.5cm 0 0]{comparison}
  \caption{Usporedba učenika koji su koristili aplikaciju s onima koji nisu}
\end{figure}

\chapter{Rasprava}

Tijekom izrade ove aplikacije uočili smo mnogo načina na koje bismo mogli
poboljšati aplikaciju. Već radimo i na novim promjenama s namjerom poboljšanja
projekta, a koje su proizašle iz dobivenih rezultata rješavanja kvizova i
ispunjavanja anketa. Jedna od dramatičnijih promjena koju smo odlučili napraviti
je učiniti aplikaciju dostupnu svima, ne samo učenicima i profesorima, gdje bi
svi mogli izrađivati kvizove i riješavati ih. Takvo bi otvorenje zahtijevalo
oprez i moglo bi promijeniti karakter aplikacije, ali mislimo da bi s tom
promijenom bilo puno zanimljivije koristiti aplikaciju.

Kako bi kvizovi dobili malo više osobnosti, bit će moguće pridružiti im slike.

Odlučili smo ukloniti vodič kroz aplikaciju jer ju želimo napraviti dovoljno
jednostavnom da upute nisu potrebne, već da je jasno samo po sebi kako postići
određeni cilj.

\chapter{Zaključak}

\chapter{Zahvale}

Htjeli bismo zahvaliti našoj mentorici, prof. dr. sc. Kristini Kocijan, za
veliku podršku u ovom projektu i prikupljanju uzorka škola koje su voljne
koristiti našu aplikaciju i pomoći nam kod istraživanja. Nadalje, htjeli bismo
zahvaliti profesorima i učenicima za sudjelovanje u ovom istraživanju i za
slanje prijedloga i prijavu grešaka, što nam je puno pomoglo kod unaprijeđivanja
aplikacije.

Sortirane prema broju aktivnih kvizova, škole koje su sudjelovale u istraživanju
su: I. gimnazija, Osnovna škola Bogumila Tonija, Osnovna škola Stjepana Radića,
Srednja škola Metković, Osnovna škola Poreč, Osnovna škola Stubičke Toplice,
Gimnazija Sesvete, Poljoprivredna škola, Srednja škola ``Ivan Seljanec''
Križevci, Gimnazija fra Dominika Mandica, Medicinska škola Varaždin, Srednja
škola Čakovec, Tehnička škola Šibenik, Privatna srednja ekonomska škola INOVA,
OŠ KV, Škola za umjetnost, dizaj, grafiku i odjeću, Zabok, Gimnazija i strukovna
škola Jurja Dobrile, Nadbiskupska klasična gimnazija.

\listoffigures

\renewcommand{\bibname}{Popis literature}

\begin{thebibliography}{9}
  \bibitem{lamport94} Leslie Lamport, \emph{\LaTeX: A Document Preparation
    System}. Addison Wesley, Massachusetts, 2nd Edition, 1994.
\end{thebibliography}

\chapter{Sažetak}

\textbf{Ključne riječi}:

\chapter{Summary}

\textbf{Keywords}:

\end{document}
